\thispagestyle{otFront}
\phantomsection\pdfbookmark[0]{Parcours de travail}{bm-parcours}
\otPageTitre{Parcours de travail}{Des bases aux résolutions de problèmes}
\fontsize{11}{15}\selectfont
\otIntertitre{Un livre de Terminale générale, spécialité physique-chimie}
Les chapitres 1 à 12 traitent de la constitution et des transformations de la matière,
les chapitres 13 à 17 du mouvement et des interactions, les chapitres 18 à 20 de l'énergie,
et les chapitres 21 à 26 des ondes et des signaux. Le chapitre 27, consacré aux mesures et
aux incertitudes, sert dès les premières séances de travaux pratiques : il peut être travaillé
à tout moment de l'année. L'ordre des chapitres peut suivre la progression du professeur.\par\vspace{5mm}
\otIntertitre{Choisir le bon point de départ}
\begin{tabularx}{\linewidth}{@{}p{40mm}Y@{}}\toprule
\textbf{Parcours} & \textbf{Objectif}\\\midrule
Découverte, 1 à 5 & Vérifier les connaissances avec des applications directes.\\[1.5mm]
Entraînement, 6 à 11 & Installer les méthodes, puis choisir une démarche.\\[1.5mm]
Objectif bac, 12 à 15 & Exploiter des documents, modéliser, rédiger une argumentation.\\[1.5mm]
Défis D01 à D\otNbDefis & Prendre des initiatives dans des situations exigeantes.\\[1.5mm]
Sujets type bac & S'entraîner en temps limité : 3 h 30, trois exercices.\\\bottomrule
\end{tabularx}
\vspace{5mm}
\otIntertitre{Une méthode de travail en quatre temps}
\begin{questions}
\item Lire les objectifs et L'essentiel. Reformuler chaque loi avec ses conditions de validité.
\item Chercher sans corrigé : schématiser la situation, écrire les données et les relations utiles.
\item Comparer sa démarche à la correction. Repérer précisément l'étape qui manquait.
\item Refaire quelques jours plus tard, puis expliquer la solution à voix haute.
\end{questions}
\vspace{3mm}
\otIntertitre{L'épreuve du baccalauréat, session 2027}
La partie écrite dure 3 h 30 et comporte trois exercices indépendants ; deux points sur vingt
y sont consacrés à la maîtrise de la langue : orthographe, syntaxe, qualité du raisonnement
et précision du vocabulaire. La partie pratique, l'évaluation des compétences expérimentales
(ECE), dure 1 heure. La note de l'épreuve vaut 0,8 fois celle de l'écrit plus 0,2 fois celle de l'ECE.
\editionRetenir{Un résultat sans unité, ou avec trop de chiffres significatifs, est incomplet.
Une phrase de conclusion qui répond à la question posée fait partie de la réponse.}
\clearpage
